% コンパイル（slide/20261007/ で実行。PDFは slide/out/20261007/ 、他は build/）:
%   ../lualatex-dated.sh -synctex=1 -interaction=nonstopmode -file-line-error -output-directory=out survey_kawamotostyle_20261007_handout_v2.tex
\documentclass[aspectratio=169,10pt,t,handout]{beamer}
\usepackage[ja]{labslide}
% 本文・見出し・脚注など、指定されている文字サイズを一律に 6pt 大きくする．
\makeatletter
\newdimen\labvtwo@size
\newskip\labvtwo@skip
\let\labvtwo@set@fontsize\set@fontsize
\def\set@fontsize#1#2#3{%
  \@defaultunits\labvtwo@size#2pt\relax\@nnil
  \advance\labvtwo@size by 1pt\relax
  \@defaultunits\labvtwo@skip#3pt\relax\@nnil
  \advance\labvtwo@skip by 1pt\relax
  \begingroup
    \edef\x{\endgroup
      \noexpand\labvtwo@set@fontsize{\unexpanded{#1}}{\strip@pt\labvtwo@size}{\the\labvtwo@skip}%
    }\x
}
\makeatother
\usetikzlibrary{fit,backgrounds,calc}

% 題名を，ページ上端と下の横棒のあいだの中央に置く．
\makeatletter
\setlength{\beamer@frametopskip}{0mm}
\makeatother
\setbeamertemplate{frametitle}{%
  \vspace{4.2mm}%
  {\usebeamerfont{frametitle}\color{ink}\insertframetitle\par}%
  \vspace{0.4mm}%
  {\color{ruleblue}\rule{\textwidth}{0.7pt}\par}%
  \vspace{1.1mm}%
  \ifx\insertframesubtitle\@empty\else
    {\footnotesize\color{muted}\insertframesubtitle\par}%
    \vspace{1.2mm}%
  \fi
}

\definecolor{cviolet}{HTML}{F3E8F6}
\definecolor{cviolettt}{HTML}{7A458E}
\renewcommand{\usecardcolor}[1]{%
  \def\cardfill{cgreen}%
  \def\cardtitlec{cgreentt}%
  \ifstrequal{#1}{red}{%
    \def\cardfill{corange}%
    \def\cardtitlec{corangett}%
  }{}%
  \ifstrequal{#1}{orange}{%
    \def\cardfill{corange}%
    \def\cardtitlec{corangett}%
  }{}%
  \ifstrequal{#1}{blue}{%
    \def\cardfill{cblue}%
    \def\cardtitlec{cbluett}%
  }{}%
  \ifstrequal{#1}{green}{%
    \def\cardfill{cgreen}%
    \def\cardtitlec{cgreentt}%
  }{}%
  \ifstrequal{#1}{violet}{%
    \def\cardfill{cviolet}%
    \def\cardtitlec{cviolettt}%
  }{}%
  \ifstrequal{#1}{plain}{%
    \def\cardfill{cplain}%
    \def\cardtitlec{cplaintt}%
  }{}%
  \ifstrequal{#1}{gray}{%
    \def\cardfill{cplain}%
    \def\cardtitlec{cplaintt}%
  }{}%
}

\graphicspath{{fig/survey/}}

\kicker{サーベイ}
\title{歩きやすい道を選択する\\四足ロボットの局所プランナ}
\date{2026年10月}

% カード内の箇条書き．同じ段のカードは第2オプションで高さをそろえる．
\newcommand{\li}{%
  \ifvmode\else\par\fi
  \noindent\hangindent=1.15em\hangafter=1
  \makebox[1.15em][l]{\textcolor{ruleblue}{\textbullet}}\ignorespaces}
\newcommand{\lii}{%
  \ifvmode\else\par\fi
  \noindent\hspace{1.0em}\hangindent=2.15em\hangafter=1
  \makebox[1.05em][l]{\textcolor{muted}{\textendash}}\ignorespaces}

\RenewDocumentCommand{\infocard}{O{green} o m m}{%
  \usecardcolor{#1}%
  \IfValueTF{#2}{\def\cardminh{#2}}{\def\cardminh{0pt}}%
  \begin{tikzpicture}[baseline=(card.north)]
    \node[
      anchor=north west,
      fill=\cardfill,
      text width=\dimexpr\linewidth-14pt\relax,
      minimum height=\cardminh,
      inner xsep=6.5pt,
      inner ysep=4pt,
      align=flush left
    ] (card) {%
      {\color{\cardtitlec}\bfseries #3}\par\vspace{0.15em}%
      {\color{body}\footnotesize #4\par}%
    };
  \end{tikzpicture}%
}

\begin{document}

\begin{frame}[plain]
  \labtitlepage{
  }
\end{frame}

\begin{frame}{研究テーマ ｜ ロボットの制御系}
  \centering
  \begin{tikzpicture}[node distance=6.5mm]
    \node[kworange, text width=24mm, minimum height=8mm] (c) {コントローラ};
    \node[kwgreen, text width=24mm, minimum height=8mm, right=of c] (p) {局所プランナ};
    \node[kwblue, text width=24mm, minimum height=8mm, right=of p] (g) {歩行モデル};
    \node[kwpurple, text width=24mm, minimum height=8mm, right=of g] (m) {モータ};
    \draw[kwarr] (c) -- (p)
      node[midway, below=4mm, font=\scriptsize, text=muted] {目標座標};
    \draw[kwarr] (p) -- (g)
      node[midway, below=4mm, font=\scriptsize, text=muted] {目標速度};
    \draw[kwarr] (g) -- (m)
      node[midway, below=4mm, font=\scriptsize, text=muted] {トルク};
  \end{tikzpicture}

  \vspace{1.5em}
  \includegraphics[width=\textwidth,height=35mm,keepaspectratio]{img001.png}

  \takeaway{一般的に、入力は3つの変換を経てモータに届く}
\end{frame}

\begin{frame}{研究テーマ ｜ 局所プランナ}
  \centering
  \begin{tikzpicture}[node distance=6.5mm]
    \node[kworange, text width=24mm, minimum height=8mm, opacity=0.2] (c) {コントローラ};
    \node[kwgreen, text width=24mm, minimum height=8mm, right=of c] (p) {局所プランナ};
    \node[kwblue, text width=24mm, minimum height=8mm, right=of p] (g) {歩行モデル};
    \node[kwpurple, text width=24mm, minimum height=8mm, right=of g, opacity=0.2] (m) {モータ};
    \draw[kwarr] (c) -- (p)
      node[midway, below=4mm, font=\scriptsize, text=muted] {目標座標};
    \draw[kwarr] (p) -- (g)
      node[midway, below=4mm, font=\scriptsize, text=muted] {目標速度};
    \draw[kwarr] (g) -- (m)
      node[midway, below=4mm, font=\scriptsize, text=muted] {トルク};
  \end{tikzpicture}

  % img2
  \includegraphics[width=\textwidth,height=35mm,keepaspectratio]{img002.png}

  \takeaway{目標速度を、「安全保障」と「好み選択」の両方に対応させる局所プランナの研究を行う}
\end{frame}

\begin{frame}{好み選択に関する論文}
  \infocard[green][22mm]{歩行コストマップからのナビゲーション}{%
    \li Where Should I Walk? Predicting Terrain Properties From Images Via Self-Supervised Learning
    \li IEEE RA-L，2019
    \li Cited by 268}
  \vspace{0.4em}
  \infocard[blue][22mm]{RGBDと強化学習を使ったナビゲーション}{%
    \li Traversability-Aware Legged Navigation by Learning from Real-World Visual Data
    \li IEEE Transactions on Robotics
    \li Cited by 13}

  \takeaway{1本目は好みをマップで表現、2本目は好みをナビゲーションで表現する}
\end{frame}

% ---------- Wellhausen ----------

\begin{frame}{Where Should I Walk?}
  \source{Wellhausen et al., IEEE RA-L, 2019．Cited by 268}
  \centering
  \includegraphics[width=\textwidth,height=52mm,keepaspectratio]{wsiw_row.png}

  \takeaway{入力 RGB からコストマップを出し，ダイクストラで経路を引く．}
\end{frame}

\begin{frame}{背景・目的}
  \source{Wellhausen et al., IEEE RA-L, 2019}
  \infocard[orange][24mm]{背景}{%
    \li 自然環境の移動は，剛体だけとは限らない
    \li 平坦でも，水・砂・土は水分量で困難さが変わる
    \lii 人が定義した同一クラスでも，特性はばらつく
    \lii 人手のアノテーションは非常に困難，または不可能}
  \vspace{0.35em}
  \infocard[green][24mm]{目的}{%
    \li RGB 画像と，足の感覚（物理特性）を紐づける
    \li RGB 画像から，歩行コストを予測する}

  \takeaway{足の感覚と画像を結び，画像の各ピクセルに対応する歩行コストの予測を目指す．}
\end{frame}

\begin{frame}{貢献}
  \source{Wellhausen et al., IEEE RA-L, 2019}
  \infocard[blue][12mm]{貢献1}{\li カメラ画像と，足の接地力を結びつける}
  \vspace{0.28em}
  \infocard[orange][12mm]{貢献2}{\li 歩きながら地面特性を予測する．指標は地面反力スコア}
  \vspace{0.28em}
  \infocard[green][12mm]{貢献3}{\li 疎かつスパースなラベルで学習する CNN}
  \vspace{0.28em}
  \infocard[violet][12mm]{貢献4}{\li 直観的なナビゲーションが実現できた}

  \takeaway{足元のセンサの特徴量を地面反力スコアと定義し，そのコストで経路を決める．}
\end{frame}

\begin{frame}{提案手法}
  \source{Wellhausen et al., IEEE RA-L, 2019}
  \centering
  \begin{tikzpicture}[
      scale=0.88, transform shape,
      font=\footnotesize,
      box/.style={kwblue, text width=30mm, minimum height=12mm, inner sep=2.5pt},
      tag/.style={font=\scriptsize\bfseries, text=muted}
    ]
    \node[box] (lab) at (0,0) {足場ラベリング};
    \node[box] (weak) at (4.15,1.15) {弱教師ラベル\\{\scriptsize (Segmentation)}};
    \node[box] (self) at (4.15,-1.15) {自己教師ラベル\\{\scriptsize (Ground Reaction Score)}};
    \node[box] (cls) at (8.35,1.15) {足場種類推定器};
    \node[box] (cost) at (8.35,-1.15) {歩行コストマップ};
    \node[box] (corr) at (12.5,1.15) {結果の相関確認用};
    \node[box] (nav) at (12.5,-1.15) {ナビゲーション\\{\scriptsize (ダイクストラ)}};
    \begin{scope}[on background layer]
      \node[fill=cblue, rounded corners=1.6pt, inner sep=2.2mm, fit=(lab)] (p1) {};
      \node[fill=corange, rounded corners=1.6pt, inner sep=2.2mm, fit=(weak)(self)] (p2) {};
      \node[fill=cgreen, rounded corners=1.6pt, inner sep=2.6mm, fit=(cls)(cost)] (p3) {};
      \node[fill=cviolet, rounded corners=1.6pt, inner sep=2.2mm, fit=(nav)] (p4) {};
    \end{scope}
    \node[tag, text=cbluett, anchor=south] at (p1.north) {貢献 1};
    \node[tag, text=corangett, anchor=south] at (p2.north) {貢献 2};
    \node[tag, text=cgreentt, anchor=south] at (p3.north) {貢献 3};
    \node[tag] at ($(p3.north)+(0,-20.8mm)$) {Mean Teacher};
    \node[tag, text=cviolettt, anchor=south] at (p4.north) {貢献 4};
    \draw[kwarr] (lab.east) -- (weak.west);
    \draw[kwarr] (lab.east) -- (self.west);
    \draw[kwarr] (weak) -- (cls);
    \draw[kwarr] (self) -- (cost);
    \draw[kwarr] (cls) -- (corr);
    \draw[kwarr] (cost) -- (nav);
  \end{tikzpicture}

  \takeaway{弱教師は地面の種類の確認に使い，自己教師の歩行コストマップを経路選択に使う．}
\end{frame}

\begin{frame}{貢献 1 ｜ 足場ラベリング}
  \source{Wellhausen et al., IEEE RA-L, 2019}
  \centering
  \includegraphics[width=\textwidth,height=30mm,keepaspectratio]{img003.png}

  \takeaway{座標変換とRGB-Dカメラの座標履歴から、各画像の接地点を特定する。}
\end{frame}

\begin{frame}{弱教師学習のラベリング}
  \source{Wellhausen et al., IEEE RA-L, 2019}
  \begin{columns}[T,onlytextwidth]
    \begin{column}{0.50\textwidth}
      \infocard[orange][40mm]{人が動画を見る}{%
        \li 歩いている区間をラベル付けする
        \lii Asphalt，Path，Grass，Dirt，Sand
        \li その 1 コマの足場にラベルが付く
        \li 画像上の画素にも，同じラベルが付く}
    \end{column}
    \hfill
    \begin{column}{0.50\textwidth}
      \centering
      \includegraphics[width=\linewidth,height=45mm,keepaspectratio]{class_dots.png}\\[0.15em]
      {\scriptsize\color{muted} Grass，Path，Sand}
    \end{column}
  \end{columns}

  \takeaway{可視化用、評価用に弱教師学習を行う。}
\end{frame}

\begin{frame}{貢献 2 ｜ 自己教師の地面反力スコア}
  \source{Wellhausen et al., IEEE RA-L, 2019}
  \begin{columns}[T,onlytextwidth]
    \begin{column}{0.50\textwidth}
      {\footnotesize\color{body} ある点の Ground Reaction Score (GRS) を定義する．}
      \vspace{0.35em}
      \infocard[blue][13mm]{1}{\li 接地中のトルクセンサの時系列}
      \vspace{0.28em}
      \infocard[orange][13mm]{2}{\li ウェーブレット変換で 13860 次元に変換}
      \vspace{0.28em}
      \infocard[green][13mm]{3}{\li PCA の第 1 主成分を GRS とする}
    \end{column}
    \hfill
    \begin{column}{0.50\textwidth}
      \centering
      \includegraphics[width=\linewidth,height=48mm,keepaspectratio]{grs_dots.png}
    \end{column}
  \end{columns}

  \takeaway{画像から、素材ごとにGRSが異なっていることがわかる}
\end{frame}

\begin{frame}{貢献 2 ｜ ラベルが疎なので Mean Teacher を使う}
  \source{Wellhausen et al., IEEE RA-L, 2019}
  \vspace*{-0.4mm}
  \begin{columns}[T,onlytextwidth]
    \begin{column}{0.50\textwidth}
      \infocard[orange][22mm]{学習するには疎すぎる}{%
        \li 足元が多い．画像上部のデータが少ない
        \li 最大で 10\% しかラベルがない}
      \vspace{0.18em}
      \infocard[blue][22mm]{Mean Teacher}{%
        \li 同じ画像に，異なるデータ拡張
        \li それぞれを異なるモデルへ入力する
        \li 同じ画像なので，二つの出力は一致するべき}
    \end{column}
    \hfill
    \begin{column}{0.50\textwidth}
      \centering
      \includegraphics[width=\linewidth,height=48mm,keepaspectratio]{sparse.png}
    \end{column}
  \end{columns}

  \takeaway{足場の点だけでは疎なので，拡張した二つの出力を一致させるように学習する。}
\end{frame}

\begin{frame}{Mean Teacher}
  \source{Wellhausen et al., IEEE RA-L, 2019}
  \centering
  \begin{tikzpicture}[scale=0.66, transform shape, box/.style={minimum height=11mm, text width=26mm}]
    \node[kw, draw=cplaintt, fill=cplain, box] (img) at (0,0) {Image};
    \node[kworange, box] (s) at (5.8,1.85) {Student Model};
    \node[kwgreen, box] (t) at (5.8,-1.85) {Teacher Model};
    \node[kwblue, box] (loss) at (10.6,0) {差分\\[0.1em]{\scriptsize LOSS}};
    \draw[kwarr] (img.east) -- (s.west)
      node[pos=0.42, above=1.6mm, font=\scriptsize, text=muted] {Augmentation A};
    \draw[kwarr] (img.east) -- (t.west)
      node[pos=0.42, below=1.6mm, font=\scriptsize, text=muted] {Augmentation B};
    \draw[kwarr] (s.east) -- (loss.north west);
    \draw[kwarr] (t.east) -- (loss.south west);
    \draw[kwarr, cgreentt] (s.south) -- (t.north)
      node[midway, right=1.6mm, font=\scriptsize, text=cgreentt] {EMA};
  \end{tikzpicture}

  \vspace{0.3em}
  \begin{columns}[T,onlytextwidth]
    \begin{column}{0.32\textwidth}
      \infocard[blue][18mm]{拡張}{\li 同じ画像に異なるデータ拡張を施す \li 二つのモデルに入力する}
    \end{column}
    \hfill
    \begin{column}{0.32\textwidth}
      \infocard[orange][18mm]{Student}{\li 同じ画像なので Teacher の出力と一致すべき．\li 差分を LOSS とする}
    \end{column}
    \hfill
    \begin{column}{0.32\textwidth}
      \infocard[green][18mm]{Teacher}{\li 直接勾配降下をしない。\li Student の重みを，移動平均をとりながら更新する}
    \end{column}
  \end{columns}

  \takeaway{ラベルの無い画素も，拡張をすることでラベルが得られ、学習できるようになる}
\end{frame}

\begin{frame}{結果 ｜ 弱教師のクラス分類}
  \source{Wellhausen et al., IEEE RA-L, 2019, Table~I}
  \vspace{0.2em}
  \centering
  {\footnotesize
  \begin{tabular}{lcccccc}
    \toprule
    & asphalt & path & grass & dirt & sand & Mean \\
    \midrule
    精度 & 93.1\% & 97.0\% & 97.3\% & 85.2\% & 80.3\% & \textbf{90.6\%} \\
    \bottomrule
  \end{tabular}
  }

  \vspace{0.8em}
  \begin{columns}[T,onlytextwidth]
    \begin{column}{0.48\textwidth}
      \infocard[green][16mm]{全体}{\li 全体平均で 90.6\%の精度。}
    \end{column}
    \hfill
    \begin{column}{0.48\textwidth}
      \infocard[orange][16mm]{dirt と sand}{\li サンプル数が少なく，精度が低いのではないか}
    \end{column}
  \end{columns}

  \takeaway{非常にスパースなラベルから得られた結果としては、非常に優秀である}
\end{frame}

\begin{frame}{結果 ｜ 自己教師の回帰}
  \source{Wellhausen et al., IEEE RA-L, 2019}
  \begin{columns}[T,onlytextwidth]
    \begin{column}{0.45\textwidth}
      \centering
      \includegraphics[width=\linewidth,height=42mm,keepaspectratio]{pca.png}
    \end{column}
    \hfill
    \begin{column}{0.55\textwidth}
      \infocard[green][22mm]{素材ごとに分かれる}{%
        \li 第 1 主成分の特徴が，地面の素材ごとに違う}
      \vspace{0.4em}
      \infocard[orange][22mm]{固いほど分散が大きい}{%
        \li 固い地面ほど接地した瞬間の振動が大きい
        \li 振動が吸収し切れていないからではないか}
    \end{column}
  \end{columns}

  \takeaway{第 1 主成分は素材ごとに分かれ，固い素材ほど分散が大きい．}
\end{frame}

\begin{frame}{結果 ｜ クラスと地面反力スコア}
  \source{Wellhausen et al., IEEE RA-L, 2019}
  \centering
  \includegraphics[width=\textwidth,height=52mm,keepaspectratio]{img004.png}
  \includegraphics[width=\textwidth,height=25mm,keepaspectratio]{img005.png}

  \takeaway{地面の素材ごとにクラス分けができ、素材ごとに特徴のあるGRSが予測できている}
\end{frame}

\begin{frame}{貢献 3 ｜ ナビゲータ}
  \source{Wellhausen et al., IEEE RA-L, 2019}
  \begin{columns}[T,onlytextwidth]
    \begin{column}{0.50\textwidth}
      \vspace{2em}
      \infocard[green][30mm]{実験設定}{%
        \li Jetson TX2 で 6\,Hz
        \li 10\,cm の 2D グリッドに地面反力スコアを予測
        \li 常に 10\,m 前方へ進むよう指示する
        \li ダイクストラ法で経路を計算する}
    \end{column}
    \hfill
    \begin{column}{0.50\textwidth}
      \centering
      \includegraphics[width=\linewidth,height=52mm,keepaspectratio]{navpath.png}
    \end{column}
  \end{columns}

  \takeaway{歩きやすい「道」を選びながら、道に沿うように歩いた}
\end{frame}

\begin{frame}{まとめ ｜ Where Should I Walk?}
  \source{Wellhausen et al., IEEE RA-L, 2019}
  \infocard[green]{貢献}{%
    \li カメラ画像と足の接地力を結びつける
    \li 歩きながら地面特性を予測する
    \lii 地面反力スコア
    \li 疎かつスパースなラベルで学習する CNN
    \li 直観的なナビゲーション}
  \vspace{0.08em}
  \infocard[orange]{課題}{%
    \li 動的環境に弱い
    \lii 少し先の地面を見ている
    \lii 動く障害物の検知には向いていない}
  \vspace{0.08em}
  \infocard[blue]{展望}{%
    \li 歩行コストを歩行器へ載せる
    \lii 今回載せたのは経路計画
    \lii 歩幅や足の上げ幅にも使えるかもしれない}

  \takeaway{歩行コストは経路計画に載った．歩行器へ載せる余地が残る．}
\end{frame}

% ---------- Zhang et al. ----------

\begin{frame}{Traversability-Aware Legged Navigation}
  \source{Zhang et al., IEEE Transactions on Robotics．Cited by 13}

  \vspace{0.25em}
  \centering
  \includegraphics[width=\textwidth,height=48mm,keepaspectratio]{trav_teaser.png}

  \takeaway{RGB と Depth から、直接速度を出すRLベースの局所プランナ}
\end{frame}

\begin{frame}{背景・目的}
  \source{Zhang et al., IEEE Transactions on Robotics}
  \infocard[orange][24mm]{背景}{%
    \li 障害物を検知し，通れるかは判断できる
    \li 見た目から摩擦や弾性を予測するのは難しい
    \li 物理特性は，その地形を歩いてみないと分からない}
  \vspace{0.35em}
  \infocard[green][18mm]{目的}{%
    \li 野外で，走行可能性を考慮した四足歩行の自立ナビゲーションを，強化学習で実現する}

  \takeaway{歩いた経験から，歩きやすい道を通れるようなアクションを選ぶ局所プランナを実現する}
\end{frame}

\begin{frame}{貢献}
  \source{Zhang et al., IEEE Transactions on Robotics}
  \infocard[blue][13mm]{貢献 1}{%
    \li 不確実性を考慮した「自己経験に基づく走破性評価」}
  \vspace{0.28em}
  \infocard[green][13mm]{貢献 2}{%
    \li RGB-D を統合した，地形を避けるナビゲーション}
  \vspace{0.28em}
  \infocard[orange][13mm]{貢献 3}{%
    \li 実機での，オンオフライン学習のパイプラインによる、サンプル効率の高い学習}

  \takeaway{走破性は，自己経験と不確実性から評価する．}
\end{frame}

\begin{frame}{問題設定}
  \source{Zhang et al., IEEE Transactions on Robotics}
  \infocard[blue][12mm]{1．目標到達}{\li 最小時間でゴールへ着く}
  \vspace{0.28em}
  \infocard[orange][12mm]{2．障害物回避}{\li 動的・静的な障害物を避ける}
  \vspace{0.28em}
  \infocard[green][12mm]{3．走行コスト}{\li 自身の経験から，歩きやすい道を選ぶ}

  \takeaway{着くこと，避けること，歩きやすいことの三つを同時に満たすことを目標にする．}
\end{frame}

\begin{frame}{提案手法 ｜ 構成}
  \source{Zhang et al., IEEE Transactions on Robotics}
  \infocard[blue][16mm]{構成}{%
    \li RGBD と目的地を入力する局所プランナ $\pi^{p}_{\mathrm{rgbd}}$
    \li 盲目の歩行プランナ $\pi^{c}$
    \li 学習時の報酬＝走行可能性コストを予測する Traversability Estimator}

  \vspace{0.2em}
  \centering
  \begin{tikzpicture}[
      scale=0.78, transform shape,
      font=\scriptsize,
      inp/.style={kwblue, minimum width=20mm, minimum height=7.2mm, inner sep=1.6pt},
      mid/.style={kwblue, align=center, inner sep=2pt}
    ]
    \node[inp] (rgb) at (0,1.05) {RGB};
    \node[inp] (dep) at (0,0) {Depth};
    \node[inp] (goal) at (0,-1.05) {目標地点};
    \node[mid, minimum width=24mm, minimum height=26mm, text width=22mm] (p) at (3.15,0)
      {局所プランナ\\$\pi^{p}_{\mathrm{rgbd}}$};
    \node[mid, minimum width=24mm, minimum height=9mm, text width=22mm] (walk) at (6.55,0.62)
      {歩行プランナ\\$\pi^{c}$};
    \node[kwgreen, align=center, minimum width=26mm, minimum height=9mm, text width=24mm, inner sep=2pt] (est) at (6.55,-0.72)
      {Traversability\\Estimator};
    \node[mid, minimum width=18mm, minimum height=26mm, text width=16mm] (rob) at (9.55,0) {ロボット};
    \node[kwgreen, align=center, minimum width=46mm, minimum height=7mm, inner sep=2pt] (buf) at (4.7,-2.45)
      {学習用バッファ：$o^{p},a^{p},r^{p}$};
    \draw[kwarr] (rgb.east) -- (rgb.east -| p.west);
    \draw[kwarr] (dep.east) -- (p.west);
    \draw[kwarr] (goal.east) -- (goal.east -| p.west);
    \draw[kwarr] (p.east |- walk.west) -- (walk.west);
    \draw[kwarr] (p.east |- est.west) -- (est.west);
    \draw[kwarr] (walk.east) -- (walk.east -| rob.west);
    \draw[kwarr] (rob.west |- est.east) -- (est.east);
    \coordinate (bus) at ($(buf.north)+(0,3.2mm)$);
    \draw[cgreentt] (p.south) -- (p.south |- bus);
    \draw[cgreentt] (est.south) -- (est.south |- bus);
    \draw[cgreentt] (rob.south) -- (rob.south |- bus);
    \draw[cgreentt] (p.south |- bus) -- (rob.south |- bus);
    \draw[kwarr, cgreentt] (bus) -- (buf.north);
  \end{tikzpicture}

  \takeaway{学習時の報酬に，走行可能性コストを入れる．}
\end{frame}

\begin{frame}{提案手法 ｜ 四つの段階}
  \source{Zhang et al., IEEE Transactions on Robotics}
  \centering
  \includegraphics[width=\textwidth,height=52mm,keepaspectratio]{stages.png}

  \takeaway{歩行器，Traversability Estimator, 深度プランナ，実世界の収集，オンライン学習の順で組む．}
\end{frame}

\begin{frame}{Stage 1 ｜ 歩行器と歩行可能性 $T$}
  \source{Zhang et al., IEEE Transactions on Robotics}
  \begin{columns}[T,onlytextwidth]
    \begin{column}{0.52\textwidth}
      \infocard[blue][12mm]{歩行器 $\pi^{c}$}{%
        \li 目標速度に追従するよう，PPO で学習する}
      \vspace{0.28em}
      \infocard[green][26mm]{Estimator $T$}{%
        \li Critic の価値 $V$ を使う
        \lii 歩きやすい道ほど価値が高い
        \li 平地と今の地形の差分をコストにする
        \lii 入力コマンドの違いを無視するため
        \li $V_{\mathrm{terrain}}$ の分散が大きいときは，速度追従誤差をコストとみる}
    \end{column}
    \hfill
    \begin{column}{0.45\textwidth}
      \centering
      \includegraphics[width=\linewidth,height=26mm,keepaspectratio]{stage1.png}
    \end{column}
  \end{columns}

  \vspace{0.02em}
  \centering
  \includegraphics[width=\textwidth,height=14.6mm,keepaspectratio]{eq1.png}

  \takeaway{歩きやすい(平地に近い)ほど価値が高い。コスト予測の分散が大きい場合は、速度追従コストを使用。}
\end{frame}

\begin{frame}{Stage 2 ｜ Depth のみのプランナ}
  \source{Zhang et al., IEEE Transactions on Robotics}
  \begin{columns}[T,onlytextwidth]
    \begin{column}{0.30\textwidth}
      \infocard[plain][16mm]{入力}{%
        \li Depth
        \li 目的地}
    \end{column}
    \hfill
    \begin{column}{0.22\textwidth}
      \infocard[blue][16mm]{出力}{\li 速度の変化量}
    \end{column}
    \hfill
    \begin{column}{0.44\textwidth}
      \infocard[green][16mm]{学習}{%
        \li シミュレータ上で学習する
        \li 2a：目的地までのベクトルと移動速度の内積
        \li 2b：前向き速度．後ろが見えないため
        \li 2c：衝突ペナルティ}
    \end{column}
  \end{columns}

  \vspace{0.2em}
  \centering
  \includegraphics[width=0.42\textwidth,height=16mm,keepaspectratio]{stage2.png}

  \vspace{0.15em}
  \centering
  \includegraphics[width=\textwidth,height=16mm,keepaspectratio]{eq2.png}

  \takeaway{実世界へ出る前に、深度だけでゴールと障害物を学ぶ深度プランナを実現する。}
\end{frame}

\begin{frame}{Stage 3 ｜ 現実世界のデータを集める}
  \source{Zhang et al., IEEE Transactions on Robotics}
  \begin{columns}[T,onlytextwidth]
    \begin{column}{0.46\textwidth}
      \infocard[orange][18mm]{方法}{%
        \li Stage 2 の Depth プランナで歩く
        \li 人が速度を与えて歩く}
      \vspace{0.12em}
      \infocard[green][30mm]{収集するデータ}{%
        \li 状態 $o^{p}_{\mathrm{rgbd}}$
        \lii RGB 特徴，Depth 特徴
        \lii ロボットの状態，目標地点へのコマンド
        \li 行動 $a^{p}$：速度の変化量
        \li 報酬：走破性コストを考慮した報酬}
    \end{column}
    \hfill
    \begin{column}{0.51\textwidth}
      \centering
      \includegraphics[width=\linewidth,height=32mm,keepaspectratio]{stage3.png}\\[0.1em]
      \includegraphics[width=\linewidth,height=12mm,keepaspectratio]{eq3.png}
    \end{column}
  \end{columns}

  \takeaway{局所プランナを学習するための、状態、行動、報酬を収集する。}
\end{frame}

\begin{frame}{Stage 4 ｜ オンライン学習}
  \source{Zhang et al., IEEE Transactions on Robotics}
  {\footnotesize\color{body} 歩きながら，RGB も含むプランナを学習する．}
  \vspace{0.15em}
  \begin{columns}[T,onlytextwidth]
    \begin{column}{0.50\textwidth}
      \infocard[blue][15mm]{Soft Actor Critic}{%
        \li 過去のデータも使う
        \li 行動をランダムにし，一つのルートにしない}
      \vspace{0.18em}
      \infocard[orange][15mm]{DroQ}{%
        \li Q ネットワーク内にドロップアウト
        \li 少ないデータでの過大評価を防ぐ}
      \vspace{0.18em}
      \infocard[green][15mm]{RLPD}{%
        \li 学習に使うバッチを管理する
        \li オフラインとオンラインを半分ずつ使う}
    \end{column}
    \hfill
    \begin{column}{0.50\textwidth}
      \centering
      \includegraphics[width=\linewidth,height=48mm,keepaspectratio]{stage4.png}
    \end{column}
  \end{columns}

  \takeaway{オフライン学習や、行動のランダム化で、少ないデータでも効率よく学習する。}
\end{frame}

\begin{frame}{結果 ｜ 走破性コストの評価}
  \source{Zhang et al., IEEE Transactions on Robotics}
  {\footnotesize\color{body} 学習時の走破性コスト $T$ を，次のどれかに置き換える．}

  \vspace{0.35em}
  \begin{columns}[T,onlytextwidth]
    \begin{column}{0.32\textwidth}
      \infocard[red][12mm]{Ours}{\li Critic の価値 $V$ を使う}
    \end{column}
    \hfill
    \begin{column}{0.32\textwidth}
      \infocard[blue][12mm]{Tracking}{\li 速度追従誤差}
    \end{column}
    \hfill
    \begin{column}{0.32\textwidth}
      \infocard[orange][12mm]{SysID}{\li 摩擦推定器による摩擦}
    \end{column}
  \end{columns}

  \vspace{0.3em}
  \centering
  \includegraphics[width=0.86\textwidth,height=32mm,keepaspectratio]{bars_cost.png}

  \takeaway{Ours が，最も低いコストで歩けている．}
\end{frame}

\begin{frame}{結果 ｜ 視覚入力の重要性}
  \source{Zhang et al., IEEE Transactions on Robotics}
  \begin{columns}[T,onlytextwidth]
    \begin{column}{0.32\textwidth}
      \infocard[red][14mm]{Ours}{\li すべてにおいてよい}
    \end{column}
    \hfill
    \begin{column}{0.32\textwidth}
      \infocard[blue][14mm]{RGB のみ}{\li 衝突が多い}
    \end{column}
    \begin{column}{0.32\textwidth}
      \infocard[orange][14mm]{Depth のみ}{\li 走破性コストが高い}
    \end{column}
    \hfill
  \end{columns}

  \vspace{0.35em}
  \centering
  \includegraphics[width=0.78\textwidth,height=34mm,keepaspectratio]{bars_vision.png}

  \takeaway{Depth だけはコストが高く，RGB だけは衝突が多いため、両方必要}
\end{frame}

\begin{frame}{結果 ｜ オンライン・オフライン学習}
  \source{Zhang et al., IEEE Transactions on Robotics}
  \vspace{0.1em}
  \centering
  {\scriptsize
  \begin{tabular}{lcc}
    \toprule
    & オフライン（過去のお手本） & オンライン（探索データ） \\
    \midrule
    \textbf{Ours（RLPD）} & 通常データ 50\% & 通常データ 50\% \\
    Behavior Cloning & 模倣データ & なし \\
    Vanilla SAC & なし & 通常データ \\
    SAC with BCL & 模倣データ & 通常データ \\
    \bottomrule
  \end{tabular}
  }

  \vspace{0.25em}
  \centering
  \includegraphics[width=0.78\textwidth,height=32mm,keepaspectratio]{bars_rlpd.png}

  \takeaway{Ours は少ないデータでも効率よく学び，高い性能を持つ．}
\end{frame}

\begin{frame}{結果 ｜ 実機での検証}
  \source{Zhang et al., IEEE Transactions on Robotics}
  \centering
  \includegraphics[width=\textwidth,height=54mm,keepaspectratio]{real_cost.png}

  \takeaway{歩きやすい路面では，走破性コストが下がる．}
\end{frame}

\begin{frame}{結果 ｜ 実機でのナビゲーション}
  \source{Zhang et al., IEEE Transactions on Robotics}
  \centering
  \includegraphics[width=\textwidth,height=56mm,keepaspectratio]{nav_real.png}

  \takeaway{歩きにくい地面を避けて，ゴールへ歩いた．}
\end{frame}

\begin{frame}{まとめ ｜ Traversability-Aware}
  \source{Zhang et al., IEEE Transactions on Robotics}
  \infocard[green][18mm]{貢献}{%
    \li 不確実性を考慮した「事故経験に基づく走破性評価」
    \li RGB-D を統合した地形回避ナビゲーション
    \li 実機でのオンオフライン学習}
  \vspace{0.35em}
  \infocard[orange][18mm]{課題}{%
    \li 階段など，高さが大きく変わる環境への適応}

  \takeaway{走破性は報酬に入り，高さの大きい地形はまだ苦手である．}
\end{frame}

\end{document}
